\begin{corollary}[general plain-cusp normalization]\label{push:src:B.2.2}\label{push:B22}
For $n=2$ and $n=12$, \ref{push:src:A.2.4}--\ref{push:src:A.2.5} give all
finite-orbit smooth invariant plain-cusp refinements,
their proper flat families, normalization charts,
finite blow-up and blow-down presentations,
multiplicities and all multiple loci. \ref{push:src:A.4.7}--\ref{push:src:A.4.9}
give their complete integral topology and nearby
Euler function. These statements concern the plain
untwisted cusp. They assert no compatibility with
the finite logarithmic-transform pieces of B.1.
\end{corollary}
\Uses{\cref{push:src:A.2.4,push:src:A.2.5,push:src:A.4.7,push:src:A.4.8,push:src:A.4.9}}
\begin{proof}

The base maps $u=w^2$ and $u=r^{12}$ induce multiplication by
$2$ and $12$ on the integral parameter characteristic. Insert
$n=2$ and $n=12$ respectively into the periodic lattice action
and fan construction of \ref{push:src:A.2.4}. Its converse
recovers the invariant subdivision from each such plain-cusp
model, so the parametrization includes every stated output.

For each primitive ray $v$, the local parameter is a monomial
whose exponent on its divisor is its integral height $\ell(v)$.
Thus the central divisor is reduced exactly when every such
height equals one. Smoothness, flatness, properness, the
normalization charts and all intersection maps are those proved
in \ref{push:src:A.2.4}; \ref{push:src:A.2.5} gives the finite
blow-up and blow-down presentation. Finally substitute the
resulting orbit counts and maps into
\ref{push:src:A.4.7,push:src:A.4.8,push:src:A.4.9} for the
claimed topology. Every step is on the specified plain-cusp
charts, so the conclusion has exactly the scope in the statement.

\end{proof}
